\section*{Chapter \ref{ch:b1:erosion-transport} --- Erosion and Transport}

\begin{solution}{exo:b1:erosion-transport:1}
\qty{0.4}{g.L^{-1}} is \qty{0.4}{kg.m^{-3}}, so the flux is $500 \times 0.4 =
\qty{200}{kg.s^{-1}}$. A year of \qty{365.25}{days} lasts
\qty{3.156e7}{s}, so the yearly flux is \qty{6.31e9}{kg}, or
\qty{6.31}{Mt}.
\end{solution}

\begin{solution}{exo:b1:erosion-transport:2}
$w_s = R g d^2/(18\nu) = 1.6548 \times 9.81 \times (\num{30e-6})^2/(18 \times
\num{1.0034e-6}) = \qty{8.09e-4}{m.s^{-1}}$, about \qty{0.8}{mm} a second. Its
grain Reynolds number is $w_s d/\nu = \num{8.09e-4} \times
\num{30e-6}/\num{1.0034e-6} = 0.024$, far below 1: Stokes' law applies. It
sinks through \qty{2}{m} in $2/\num{8.09e-4} = \qty{2472}{s}$, about
\qty{41}{min}.
\end{solution}

\begin{solution}{exo:b1:erosion-transport:3}
$\tau_b = \rho g h S = 998.2 \times 9.81 \times 0.5 \times \num{4e-4} =
\qty{1.96}{Pa}$. With $(\rho_s - \rho) g = 1651.8 \times 9.81 =
\qty{16204}{N.m^{-3}}$: for \qty{1}{mm}, $\theta = 1.96/(16204 \times 0.001) =
0.121$, nearly four times the threshold of \num{0.0314}, so the sand moves; for
\qty{10}{mm}, $\theta = 0.0121$, a fifth of the threshold of \num{0.0556}, so the
gravel stays. The stream carries its sand over a gravel bed that does not move.
\end{solution}

\begin{solution}{exo:b1:erosion-transport:4}
(a)~$12.6/3.9 = 3.2$: the rivers carry about three times more in suspension
than in solution. (b)~Large reservoirs: $0.20 \times 14 = \qty{2.8}{Gt}$ a
year; all of them: $0.26 \times 14 = \qty{3.64}{Gt}$ a year. (c)~The dissolved
flux is \qty{3.9e15}{g} a year at \qty{120}{g.m^{-3}}, so the water carrying it
is $\num{3.9e15}/120 = \qty{3.25e13}{m^3}$, about \qty{32500}{km^3} a year. The
\omterm{def:b1:erosion-transport:load}{suspended load} in that water is $\qty{1.26e16}{g}/\qty{3.25e16}{L} =
\qty{0.39}{g.L^{-1}}$.
\end{solution}

\begin{solution}{exo:b1:erosion-transport:5}
(a)~The section has the area $10 \times 2 = \qty{20}{m^2}$ and the wetted
perimeter $10 + 2 \times 2 = \qty{14}{m}$, so $R_h = 20/14 = \qty{1.43}{m}$.
(b)~$\tau_b = 998.2 \times 9.81 \times 1.43 \times 0.001 = \qty{14.0}{Pa}$,
against $\qty{19.6}{Pa}$ for a very wide channel \qty{2}{m} deep: the ratio is
$R_h/h = 0.714$, the banks taking the rest of the friction. (c)~For a
rectangular channel of width $b$, $R_h = bh/(b + 2h) = h/(1 + 2h/b)$, which
tends to $h$ when $b$ grows.
\end{solution}

\begin{solution}{exo:b1:erosion-transport:6}
(a)~$d_1 = \left(18\nu^2/(Rg)\right)^{1/3} = (18/(Rg))^{1/3}\,\nu^{2/3}$, and $R$
and $g$ do not depend on the viscosity. (b)~$\nu = \num{1.3059e-3}/999.7 =
\qty{1.306e-6}{m^2.s^{-1}}$ at \qty{10}{\celsius}, against
\qty{1.0034e-6}{m^2.s^{-1}} at \qty{20}{\celsius}: the ratio is \num{1.302},
and $1.302^{2/3} = 1.192$, so $d_1 = 0.104 \times 1.192 = \qty{0.124}{mm}$.
(c)~In the Stokes range $w_s$ is inversely proportional to $\nu$: the grain
settles $1.302$ times more slowly, about \qty{30}{\percent}. Cold water holds
fine sediment in suspension longer: a river in winter carries its silt
somewhat farther for the same flow.
\end{solution}

\begin{solution}{exo:b1:erosion-transport:7}
With $R g/(18\nu) = \qty{8.99e5}{m^{-1}.s^{-1}}$ and the combined law:
\begin{center}
\begin{tabular}{@{}lccc@{}}
\hline
$d$ & Stokes (\si{m.s^{-1}}) & combined (\si{m.s^{-1}}) & ratio\\
\hline
\qty{0.1}{mm} & \num{8.99e-3} & \num{7.53e-3} & \num{1.19}\\
\qty{0.5}{mm} & \num{0.225} & \num{0.0711} & \num{3.16}\\
\qty{2}{mm} & \num{3.60} & \num{0.197} & \num{18.3}\\
\hline
\end{tabular}
\end{center}
The ratio grows with the size, because the grain Reynolds number grows: Stokes'
law, valid only for small Reynolds numbers, increases as $d^2$, while the drag of
the wake slows the coarse grains to a velocity that grows only as $\sqrt d$.
\end{solution}

\begin{solution}{exo:b1:erosion-transport:8}
(a)~$(Rg/\nu^2)^{1/3} = \qty{25263}{m^{-1}}$. At \qty{0.2}{mm}, $D_* = 5.05$
and $\theta_c = 0.3/7.06 + 0.055(1 - e^{-0.101}) = 0.0425 + 0.0053 =
0.0478$; $\tau_c = 0.0478 \times 16204 \times \num{2e-4} = \qty{0.155}{Pa}$.
At \qty{20}{mm}, $D_* = 505$ and $\theta_c = 0.3/607 + 0.055(1 - e^{-10.1}) =
0.0005 + 0.0550 = 0.0555$; $\tau_c = 0.0555 \times 16204 \times 0.02 =
\qty{18.0}{Pa}$. (b)~$h = \tau_c/(\rho g S)$ with $\rho g S =
\qty{19.58}{Pa.m^{-1}}$: \qty{7.9}{mm} of water moves the fine sand, \qty{0.92}{m}
the gravel.
\end{solution}

\begin{solution}{exo:b1:erosion-transport:9}
The basin loses \qty{0.8}{Mt} = \qty{8e5}{t} a year, or $\num{8e5}/2000 =
\qty{400}{t.km^{-2}.yr^{-1}}$. By \cref{prop:b1:erosion-transport:denudation},
$D = \qty{8e8}{kg}/(2650 \times \qty{2e9}{m^2}) = \qty{1.51e-4}{m}$ a year, or
\qty{151}{m/Myr}. That is about three times the median of the basins measured
with beryllium-10, \qty{54}{m/Myr}, though below their mean of
\qty{218}{m/Myr}: a fairly fast basin, as in steep or tectonically active
country.
\end{solution}

\begin{solution}{exo:b1:erosion-transport:10}
(a)~$u_* = \sqrt{g h S} = \sqrt{9.81 \times 3 \times \num{5e-4}} =
\qty{0.1213}{m.s^{-1}}$ and $\tau_b = \rho u_*^2 = \qty{14.7}{Pa}$.
(b)~Suspension needs $w_s < u_* = \qty{0.1213}{m.s^{-1}}$. The combined law
gives \qty{0.126}{m.s^{-1}} at \qty{1}{mm}; solving
$w_s(d) = 0.1213$ by trial gives $d = \qty{0.94}{mm}$. (c)~For gravel
$\theta_c \approx 0.0555$, so $\tau_c = \qty{14.7}{Pa}$ for $d \approx
14.7/(0.0555 \times 16204) = \qty{16.3}{mm}$; at that size the threshold curve gives
$\theta_c = 0.0555$ indeed. (d)~Grains coarser than about \qty{16}{mm} stay on
the bed and form a stable gravel floor; finer gravel and the sand coarser than
about a millimetre travel as \omterm{def:b1:erosion-transport:load}{bed load}; finer sand, silt and clay travel in
suspension.
\end{solution}

\begin{solution}{exo:b1:erosion-transport:11}
(a)~$0.30 \times \qty{1000}{m} = \qty{300}{m}$ of rock, averaged over the
region. (b)~The rebound is $300 \times 2700/3300 = \qty{245}{m}$. (c)~The mean
surface is lowered by $300 - 245 = \qty{55}{m}$ only. (d)~The summits, which
were not eroded, are carried up with the whole region: they now stand
\qty{245}{m} higher than before. Cutting valleys can raise the peaks of a range.
\end{solution}

\begin{solution}{exo:b1:erosion-transport:12}
(a)~$0.9 \times \qty{5e8}{kg}/\qty{1300}{kg.m^{-3}} = \qty{346000}{m^3}$ a year,
or \qty{0.69}{\percent} of the reservoir. (b)~Half the reservoir,
\qty{25e6}{m^3}, is filled in $\num{25e6}/\num{346000} = \qty{72}{years}$ at this
rate. (c)~As the reservoir fills, its volume, and so the residence time of the
water, decreases, and finer grains pass through: its trapping efficiency falls.
Meanwhile the older deposits are buried and compacted, expelling water, so that
their bulk density rises and each tonne takes less room. Both effects slow the
loss of volume as the reservoir ages.
\end{solution}

\begin{solution}{pb:b1:erosion-transport:1}
\textbf{1.} A year of water is $1.2 \times \num{3.156e7} = \qty{3.79e7}{m^3}$;
spread over \qty{8e7}{m^2}, a layer of \qty{0.473}{m}: a runoff of
\qty{473}{mm} a year.

\textbf{2.} $1.2 \times \qty{0.150}{kg.m^{-3}} = \qty{0.18}{kg.s^{-1}}$, or
\qty{5680}{t} a year.

\textbf{3.} $1.2 \times 0.080 = \qty{0.096}{kg.s^{-1}}$, or \qty{3030}{t} a
year.

\textbf{4.} The \omterm{def:b1:erosion-transport:load}{bed load} is \qty{568}{t} a year, and the total $5680 + 3030 +
568 = \qty{9278}{t}$ a year.

\textbf{5.} The concentration of suspended sediment is high during the few days
of flood, when the discharge is high too, and low the rest of the time. A
time average gives every day the same weight, so it is dominated by the many
days of low flow and low concentration. By
\cref{prop:b1:erosion-transport:flux}, the yearly mass is
$\int Q\,C\,\mathrm dt$; the flow-weighted mean is this integral divided by
$\int Q\,\mathrm dt$, so it is the flow-weighted mean that must be multiplied by
the mean discharge.

\textbf{6.} $3030/9278 = \num{0.33}$, a third. For the world's rivers,
$3.9/(12.6 + 3.9) = \num{0.24}$, and less if the \omterm{def:b1:erosion-transport:load}{bed load} were counted.
The basin's share is higher: low hills, where slow mechanical erosion leaves
the rock time to dissolve, and limestone and chalk, which dissolve
completely.

\textbf{7.} $9278/80 = \qty{116}{t.km^{-2}.yr^{-1}}$.

\textbf{8.} $D = \qty{9.278e6}{kg}/(2700 \times \qty{8e7}{m^2}) =
\qty{4.30e-5}{m}$ a year, \qty{43.0}{m/Myr}.

\textbf{9.} $43.0/54 = 0.80$: a little below the median of the measured
basins, and five times below their mean. A moderately slow basin, as expected
for low hills in a quiet region.

\textbf{10.} The gauging averages over the few years of measurement, and depends
on the floods that happened in them. Beryllium-10 averages over the time taken
to erode the few metres in which it is produced: at \qty{43}{m/Myr}, each metre
takes \num{23000} years, so the average covers tens of thousands of years. They
can differ because of rare large floods, missing from a short record, because
farming and clearing have raised erosion in recent centuries, and because
sediment is stored in, and taken back from, the floodplain.

\textbf{11.} Part of the \omterm{def:b1:erosion-transport:load}{dissolved load} does not come from the rock: much of the
bicarbonate carries carbon taken from the air or the soil, and rain brings sea
salt. Counting them would credit the basin with rock it has not lost.

\textbf{12.} $\num{1.5e6}/1.2 = \qty{1.25e6}{s}$, \num{14.5} days.

\textbf{13.} The combined law gives \qty{7.53e-3}{m.s^{-1}}: the grain falls
\qty{8}{m} in \qty{1062}{s}, about \qty{18}{min}. It is trapped.

\textbf{14.} $w_s = \num{8.99e5} \times (\num{2e-5})^2 =
\qty{3.60e-4}{m.s^{-1}}$: \qty{8}{m} in \qty{22200}{s}, about \num{6.2}~hours.
Trapped.

\textbf{15.} $w_s = \num{8.99e5} \times (\num{1e-6})^2 =
\qty{8.99e-7}{m.s^{-1}}$: \qty{8}{m} takes \qty{8.9e6}{s}, \num{103}~days, seven
times the residence time. Most of it passes through the dam with the water.

\textbf{16.} The grain must fall \qty{8}{m} in \qty{1.25e6}{s}, at
$\qty{6.4e-6}{m.s^{-1}}$ at least. By Stokes' law $d = \sqrt{18\nu
w_s/(Rg)} = \sqrt{\num{6.4e-6}/\num{8.99e5}} = \qty{2.7}{\micro m}$: a fine
clay. Coarser grains are trapped wherever they enter.

\textbf{17.} $568 + 0.7 \times 5680 = \qty{4544}{t}$ a year, which fills
$\num{4.544e6}/1300 = \qty{3496}{m^3}$ a year, \qty{0.233}{\percent} of the
reservoir. Half of it, \qty{7.5e5}{m^3}, is filled in about \num{215} years.

\textbf{18.} $100/43.0 = \qty{2.33}{Myr}$.

\textbf{19.} $1 - 2700/3300 = 0.18$: the surface comes down by only
\qty{18}{\percent} of the thickness removed.

\textbf{20.} $100/0.1818 = \qty{550}{m}$ of rock.

\textbf{21.} The whole region rises by $550 \times 2700/3300 = \qty{450}{m}$
of rebound. Summits that are hardly eroded are carried up with it, by up to
\qty{450}{m}, while the valleys deepen: the relief grows before it declines.

\textbf{22.} The \omterm{def:b1:erosion-transport:denudation}{denudation rate} depends above all on the slopes. As the hills
are lowered and their slopes ease, the rate will decrease, unless uplift or a
change of climate steepens them again: the time computed at today's rate is a
lower bound.

\textbf{23.} $550/42.95 = 12.80$. At today's rate and with isostatic rebound,
lowering the Varden Hills by \qty{100}{m} takes \textbf{12.8}~million years.
\end{solution}
